\documentclass[11pt]{beamer}
\usepackage[utf8]{inputenc}
\usepackage{amsmath, amssymb, graphicx}
\usepackage{bm}

% Presentation Theme
\usetheme{Madrid}
\usecolortheme{default}

\title[The B-Plane]{Astrodynamics 2026-1}
\subtitle{Interplanetary Targeting: \\ Understanding the B-Plane Coordinate System}
\author{Prof. Juan}
\institute{Universidad de Antioquia - Aerospace Engineering}
\date{Spring 2026}

\begin{document}

% --- Title Slide ---
\begin{frame}
    \titlepage
\end{frame}

% ==========================================
% SECTION 1: WHAT IS THE B-PLANE?
% ==========================================
\section{Defining the Target}

% --- Slide 1 ---
\begin{frame}{The Targeting Dilemma}
    When approaching a planet from millions of kilometers away, targeting a specific Keplerian element (like periapsis radius or inclination) is numerically difficult. The coordinate space is highly non-linear.
    
    \vspace{0.3cm}
    \textbf{The Solution: The B-Plane}
    Instead of targeting the orbit directly, mission designers target a specific $(x, y)$ coordinate on an imaginary 2D plane located next to the target planet. 
    
    \vspace{0.2cm}
    By changing where we "pierce" this plane, we uniquely define the resulting hyperbolic flyby or orbit insertion.
\end{frame}

% --- Slide 2 ---
\begin{frame}{Defining the B-Plane}
    The \textbf{B-Plane} is a plane passing through the center of the target body, oriented strictly perpendicular to the incoming hyperbolic excess velocity vector ($\bm{v}_\infty$).
    
    \vspace{0.3cm}
    \textbf{The B-Vector ($\bm{B}$):}
    \begin{itemize}
        \item The vector $\bm{B}$ lies entirely within the B-Plane.
        \item It points from the center of the planet to the point where the incoming asymptote pierces the plane.
        \item The magnitude of this vector, $|\bm{B}|$, is the \textbf{Impact Parameter} (the theoretical closest approach distance if the planet had no gravity).
    \end{itemize}
\end{frame}

% ==========================================
% SECTION 2: THE S, T, R COORDINATE SYSTEM
% ==========================================
\section{The S, T, R Frame}

% --- Slide 3 ---
\begin{frame}{The $\bm{S}, \bm{T}, \bm{R}$ Coordinate System}
    To define coordinates on this plane, we construct an orthogonal basis triad centered on the target planet:
    
    \begin{enumerate}
        \item \textbf{The $\bm{S}$ Vector:} Points parallel to the incoming velocity ($\bm{v}_\infty$). It is normal to the B-Plane.
        $$ \bm{S} = \frac{\bm{v}_\infty}{|\bm{v}_\infty|} $$
        \item \textbf{The $\bm{T}$ Vector:} Lies in the B-Plane. It is typically defined to be parallel to the ecliptic plane (or the planet's equatorial plane). It points roughly "horizontal".
        \item \textbf{The $\bm{R}$ Vector:} Completes the right-handed set. It points generally "vertical".
        $$ \bm{R} = \bm{S} \times \bm{T} $$
    \end{enumerate}
\end{frame}

% --- Slide 4 ---
\begin{frame}{Projecting the B-Vector}
    Since the B-vector lies entirely in the B-Plane, it can be expressed as a 2D coordinate using the $\bm{T}$ and $\bm{R}$ axes. We use the dot product to find these projections.
    
    \vspace{0.3cm}
    \textbf{The Target Coordinates:}
    \begin{itemize}
        \item \textbf{$\bm{B \cdot T}$:} The horizontal offset. In GMAT, this is \texttt{BdotT}.
        \item \textbf{$\bm{B \cdot R}$:} The vertical offset. In GMAT, this is \texttt{BdotR}.
    \end{itemize}
    
    \vspace{0.3cm}
    The total impact parameter (the aiming radius) is simply the hypotenuse:
    $$ |\bm{B}| = \sqrt{(\bm{B \cdot T})^2 + (\bm{B \cdot R})^2} $$
\end{frame}

% ==========================================
% SECTION 3: MISSION DESIGN IMPLICATIONS
% ==========================================
\section{Flying the Mission}

% --- Slide 5 ---
\begin{frame}{How the B-Plane Controls the Trajectory}
    The specific coordinates $(\bm{B \cdot T}, \bm{B \cdot R})$ dictate the exact 3D geometry of the flyby:
    
    \begin{itemize}
        \item \textbf{Closest Approach ($r_p$):} Controlled by the magnitude $|\bm{B}|$. Moving the target point outward (larger B) results in a higher altitude flyby and a smaller Turn Angle ($\delta$).
        \item \textbf{Orbital Plane (Inclination):} Controlled by the angle $\theta = \tan^{-1}\left(\frac{\bm{B \cdot R}}{\bm{B \cdot T}}\right)$. 
    \end{itemize}
    
    \vspace{0.3cm}
    \textit{Example:} To enter a perfectly equatorial orbit, you aim purely along the $\bm{T}$ axis ($\bm{B \cdot R} = 0$). To enter a polar orbit (like our Jupiter mission), you aim offset along the $\bm{R}$ axis!
\end{frame}

% --- Slide 6 ---
\begin{frame}{Gravity Assists and the B-Plane}
    For a gravity assist, the B-Plane target is the ultimate control knob. 
    
    \vspace{0.3cm}
    By tweaking \texttt{BdotT} and \texttt{BdotR} in a Differential Corrector:
    \begin{enumerate}
        \item We dictate the exact Turn Angle $\delta$ (via the impact parameter magnitude).
        \item We dictate the 3D plane in which that turn happens.
        \item We aim the outbound $\bm{v}_\infty$ vector precisely toward the next planet!
    \end{enumerate}
    
    \vspace{0.2cm}
    \textbf{Summary:} The B-Plane turns 3D orbital mechanics into a simple game of throwing darts at a 2D board.
\end{frame}

\end{document}